\hojaejercicio{Ceros y singularidades de funciones complejas}
%ejercicio1
\ejercicio{Sea $D=D(0;1)$ y $f \in H(D)$ tal que $|f(z)| \le 1$. Usa el Lema de Schwarz para probar que:
    \[
    \frac{|f'(z)|}{1-|f(z)|^2} \le \frac{1}{1-|z|^2} \quad \forall z \in D
    \]
    y que la igualdad se cumple si y sólo si $f$ es una transformación de Möbius tal que $f(D)=D$.}{
        Para un punto $a \in D$, definimos la transformación de Möbius:
\[ \phi_a(\zeta) = \frac{\zeta - a}{1 - \bar{a}\zeta} \]
Recordamos que $\phi_a$ es un automorfismo del disco unidad y su derivada en el punto $a$ es:
\[ \phi_a'(a) = \frac{1}{1 - |a|^2} \]
Asimismo, la derivada de su inversa $\phi_a^{-1}$ (que es $\phi_{-a}$) evaluada en el origen es:
\[ (\phi_a^{-1})'(0) = 1 - |a|^2 \]


Fijamos un punto arbitrario $z \in D$. Definimos la función auxiliar $g: D \to D$ como la composición:
\[ g = \phi_{f(z)} \circ f \circ \phi_z^{-1} \]
Esta función cumple las siguientes propiedades:
\begin{itemize}
    \item Es holomorfa por ser composición de funciones holomorfas.
    \item Envía el disco al disco ($|g(w)| \leq 1$ para $w \in D$).
    \item  $g(0) = \phi_{f(z)}(f(\phi_z^{-1}(0))) = \phi_{f(z)}(f(z)) = 0$.
\end{itemize}



Puesto que $g(0) = 0$ y $|g(w)| \leq 1$, por el Lema de Schwarz estándar tenemos que $|g'(0)| \leq 1$. Calculamos $g'(0)$ mediante la regla de la cadena:
\[ g'(0) = \phi_{f(z)}'(f(z)) \cdot f'(z) \cdot (\phi_z^{-1})'(0) \]
Sustituyendo los valores de las derivadas calculadas en la sección 1:
\[ g'(0) = \frac{1}{1 - |f(z)|^2} \cdot f'(z) \cdot (1 - |z|^2) \]
Imponiendo la condición $|g'(0)| \leq 1$:
\[ \left| \frac{f'(z)(1 - |z|^2)}{1 - |f(z)|^2} \right| \leq 1 \]
Lo cual equivale a:
\[ \frac{|f'(z)|}{1 - |f(z)|^2} \leq \frac{1}{1 - |z|^2} \]


La igualdad se da si y solo si $|g'(0)| = 1$. Según el Lema de Schwarz, esto ocurre únicamente cuando $g(w) = e^{i\theta}w$ para algún $\theta \in \mathbb{R}$. En ese caso:
\[ \phi_{f(z)} \circ f \circ \phi_z^{-1} = e^{i\theta}w \implies f = \phi_{f(z)}^{-1} \circ (e^{i\theta} \text{id}) \circ \phi_z \]
Como $f$ es una composición de transformaciones de Möbius, $f$ resulta ser una transformación de Möbius del disco en el disco.
    }

%ejercicio2
\ejercicio{Sea $\mathcal{F} = \{ f \in H(D) / |f(z)| \le 1 \ \forall z \in D \}$. Para un $z \in D$ fijo, determina $\sup_{f \in \mathcal{F}} |f'(z)|$ y prueba que ese supremo se alcanza. (Ind.: usa el ej 1).}{
    Sea $\mathcal{F} = \{ f \in H(D) \mid |f(z)| \leq 1, \, \forall z \in D \}$ la familia de funciones holómarfas acotadas en el disco unidad. Para un punto $z \in D$ fijo, determina el valor de:
\[ \sup_{f \in \mathcal{F}} |f'(z)| \]
y demuestra que dicho supremo se alcanza en $\mathcal{F}$.


Partimos del resultado demostrado en el \textbf{ejercicio 1},  para cualquier función $f \in \mathcal{F}$:
\begin{equation}
\frac{|f'(z)|}{1 - |f(z)|^2} \leq \frac{1}{1 - |z|^2} \quad \forall z \in D
\end{equation}
Para un punto $z$ fijo, podemos despejar $|f'(z)|$ multiplicando por el término $(1 - |f(z)|^2)$, que es siempre no negativo:
\begin{equation}
|f'(z)| \leq \frac{1 - |f(z)|^2}{1 - |z|^2}
\end{equation}



Observamos que en la expresión anterior, el denominador $1 - |z|^2$ es constante para el punto fijado. Por tanto, la parte derecha de la desigualdad se maximiza cuando el numerador $1 - |f(z)|^2$ es lo más grande posible. 

Dado que $|f(z)| \geq 0$, el valor máximo del numerador es $1$, el cual se alcanza si y solo si $|f(z)| = 0$. De esto se deduce que para toda $f \in \mathcal{F}$:
\[ |f'(z)| \leq \frac{1 - 0}{1 - |z|^2} = \frac{1}{1 - |z|^2} \]
Así, el valor $M = \frac{1}{1 - |z|^2}$ es una cota superior del conjunto $\{ |f'(z)| : f \in \mathcal{F} \}$.



Para demostrar que esta cota superior es realmente el supremo (y un máximo), consideramos la transformación de Möbius definida como:
\[ \phi_z(\zeta) = \frac{\zeta - z}{1 - \bar{z}\zeta} \]
Esta función $\phi_z$ pertenece a $\mathcal{F}$ por ser un automorfismo del disco. Al evaluar la función y su derivada en el punto $z$ obtenemos:
\begin{itemize}
    \item  $\phi_z(z) = \frac{z - z}{1 - |z|^2} = 0$.
    \item Como se calculó previamente, $|\phi_z'(z)| = \frac{1}{1 - |z|^2}$.
\end{itemize}


Puesto que existe una función en la familia (la transformación de Möbius $\phi_z$) tal que el valor de su derivada en $z$ iguala a la cota superior hallada, concluimos que:
\[ \sup_{f \in \mathcal{F}} |f'(z)| = \max_{f \in \mathcal{F}} |f'(z)| = \frac{1}{1 - |z|^2} \]
}

%ejercicio3
\ejercicio{Sea $f: \mathbb{C} \to \mathbb{C}$ una función entera. Demuestra que si existen $M, R > 0$ y $k \in \mathbb{N}$ tales que $|f(z)| \le M|z|^k$ $\forall z$ con $|z|>R$, entonces $f$ es un polinomio de grado $\le k$. (Ind.: usa el desarrollo en serie de potencias).}{
  Sea entonces $f \in H(\mathbb{C})$, por lo tanto sea $z_0=0$ se da que:
\[f(z) = \sum C_n z^n \quad \forall z \in D(0, \infty) = \mathbb{C}\]
siendo $C_n$ el coeficiente dado por la fórmula integral de Cauchy:
\[C_n = \frac{1}{2\pi i} \int_{C(0,r)} \frac{f(w)}{w^{n+1}} dw \quad \text{con } 0 < r < \infty\]

Por lo tanto, queremos demostrar que $C_n = 0 \quad \forall n > k$, es decir, que $f$ es un polinomio de grado $\le k$.

Para ello, observamos que:
\begin{align*}
|C_n| &= \left| \frac{1}{2\pi i} \int_{C(0,r)} \frac{f(w)}{w^{n+1}} dw \right| \\
&\le \frac{1}{2\pi} \max_{w \in C(0,r)} \left| \frac{f(w)}{w^{n+1}} \right| \cdot 2\pi r \\
&= \max_{w \in C(0,r)} |f(w)| \cdot \frac{1}{r^{n+1}} \cdot r \\
&= \frac{\max_{w \in C(0,r)} |f(w)|}{r^n}
\end{align*}

Por la condición anterior, podemos coger $r > R$ tal que $\max_{w \in C(0,r)} |f(w)| \le M |w|^k$. Sustituyendo en la desigualdad anterior:
\[ |C_n| \le \frac{M r^k}{r^n} = M r^{k-n} \]

Por lo tanto $|C_n|$ está acotado por:
\[|C_n| \le \lim_{r \to \infty} M r^{k-n} = \begin{cases} 
M & \text{si } n=k \\ 
0 & \text{si } n>k \\ 
\infty & \text{si } n<k 
\end{cases} \]
Es decir, $|C_n| \le 0 \quad \forall n > k \implies |C_n| = 0 \quad \forall n > k$.

Es decir, $f$ es de la forma:
\[ f(z) = \sum_{n=0}^{k} C_n z^n \]
lo cual es un polinomio de grado $\le k$.
}

%ejercicio4
\ejercicio{Sea $G = \{ z / |\Re (z)| < 1 \text{ y } |\Im(z)| < 1 \}$ y sea $f$ holomorfa en $G$ y continua en $\overline{G}$ tal que $f(z)=0$ si $\text{Re} z = 1$. Considerando $g(z) = f(z)f(iz)f(-z)f(-iz)$, demuestra que $f \equiv 0$ en $\overline{G}$.}{
    Sea $f: \mathbb{C} \to \mathbb{C}$. Tenemos que $f \in H(G)$ y $f$ continua en $\overline{G}$.

\vspace{0.3cm}

Definimos entonces $g(z) = f(z)f(\overline{z})f(-z)f(-\overline{z})$. ¿Dónde está definida $g$?
Si $\overline{G} = \{z / |\Re{(z)}| \le 1 \text{ y } |\Im{(z)}| \le 1\}$, entonces debemos asegurar que si $z \in \overline{G}$ entonces $\-iz, -z, iz \in \overline{G}$ para que $g$ esté definida en $\overline{G}$.

\vspace{0.3cm}

Entonces, si $z = a+bi$ con $|a| \le 1$ y $|b| \le 1$ entonces:
\[ iz = ia-b \implies \begin{cases} |\Im{(iz)}| = |-b| = |b| \le 1 \\ |\Re{(iz)}| = |a| \le 1 \end{cases} \implies iz \in \overline{G} \]

\[ -iz = -ai+b \implies \begin{cases} |\Im{(-iz)}| = |b| \le 1 \\ |\Re{(-iz)}| = |-a| = |a| \le 1 \end{cases} \implies -iz \in \overline{G} \]

Lo mismo ocurre con $-z \in \overline{G}$.

\vspace{0.3cm}

Entonces $g$ está definida en $\overline{G}$ y además, es continua en $\overline{G}$ y holomorfa en $G$ por ser multiplicación de funciones continuas en $\overline{G}$ y holomorfas en $G$. $\implies g \in H(G) \cap C(\overline{G})$.

Como $G$ es un dominio, $g \in H(G)$ y $g$ continua en $\overline{G}$ que es compacto, entonces $|g|$ alcanzará el máximo en $\partial G$.

\vspace{0.3cm}

Así pues, como $f(z)=0 \quad \forall z \in \partial G$. Sabiendo que:
\[ \partial G = [-1+i, 1+i] \cup [1+i, 1-i] \cup [1-i, -1-i] \cup [-1-i, -1+i]\]

Tenemos:
\begin{itemize}
    \item Si $z \in [-1+i, 1+i] \implies -iz \in [1-i, 1+i] \implies \Re{(-iz)} = 1 \implies f(-iz) = 0 \implies g(z)=0$
    \item Si $z \in [1-i, -1-i] \implies iz \in [1-i, 1+i] \implies \Re{(iz)} = 1 \implies f(iz) = 0 \implies g(z)=0$
    \item Si $z \in [-1-i, -1+i] \implies -z \in [1-i, 1+i] \implies \Re{(-z)} = 1 \implies f(-z) = 0 \implies g(z)=0$
\end{itemize}

Es decir: $g(z)=0 \quad \forall z \in \partial G$.

Como $|g|$ alcanza el máximo en $\partial G$ y $g=0$ en $\partial G \implies |g|=0 \quad \forall z \in \overline{G} \implies g(z)=0$.
Como $g(z)=0 \quad \forall z \in \overline{G} \implies f(z)=0$ en $\overline{G}$.
}
%ejercicio5
\ejercicio{Determina y clasifica las singularidades de las funciones siguientes:
    \begin{enumerate}
        \item $\displaystyle \frac{\cot z}{z^2}$
        \item $\displaystyle \sin \frac{1}{z}$
        \item $\displaystyle z^2 \sin z$
        \item $\displaystyle \frac{\sin z}{z^2(z+1)(z-2)^3}$
        \item $\displaystyle \frac{z}{1+z e^{1/z}}$ \quad (sólo en $z=0$)
    \end{enumerate}}{
a)
 \[
f(z) = \frac{\mathrm{cotan} z}{z^2} = \frac{\mathrm{sen} z}{\cos z \cdot z^2} \quad \forall z \in \mathbb{C}/A 
\]
siendo \( A = \{z \in \mathbb{C} : z \neq \frac{\pi}{2} + \pi k, \, k \in \mathbb{Z}\} \cup \{0\} \).

Sacamos \( A \) sabiendo que los ceros de \( \cos z \) en \( \mathbb{C} \) son los ceros de \( \cos z \) en \( \mathbb{R} \).
Entonces, \( f \in H(\mathbb{C}/A) \). Entonces, para cada punto \( z_0 \in A \) se da que \( \exists \delta_{z_0} \) tal que \( D'(z_0, \delta_{z_0}) \subseteq \mathbb{C}/A \Rightarrow f \in H(D'(z_0, \delta_{z_0})) \).
\( \forall z_0 \in \mathbb{C}/A \) podemos coger \( \delta_{z_0} = \delta = \frac{\pi}{4} \) ejemplo.
Pero lógicamente, \( f \notin H(D(z_0, \delta)) \). Entonces \( A \) es el conjunto de singularidades.

Sea \( z_0 \in A \) t.q \( z_0 = 0 \). Entonces vemos si \( \lim_{z \to 0} (z-0) f(z) = 0 \)
\[
\lim_{z \to 0} z f(z) = \lim_{z \to 0} \frac{\mathrm{sen} z}{\cos z \cdot z} = \lim_{z \to 0} \frac{\mathrm{sen} z}{z} \cdot \frac{1}{\cos z} = \mathrm{sen}'(0) \cdot 1 = \cos(0) \cdot 1 = 1 \neq 0
\]
Entonces \underline{no} puede ser una discontinuidad evitable.

Veamos:
\[
\lim_{z \to 0} z^2 f(z) = \lim_{z \to 0} \frac{\mathrm{sen} z}{\cos z} = \lim_{z \to 0} \tan z = 0
\]

\( \forall k > 2 \) se da que
\[
\lim_{z \to 0} z^k f(z) = \lim_{z \to 0} z^{k-2} z^2 f(z) = 0 \quad (\exists \lim_{z \to 0} f = L, \, \exists \lim_{z \to 0} g = \omega \Rightarrow \exists \lim_{z \to z_0} f \cdot g = L \cdot \omega)
\]
Por lo tanto, por la caracterización de polos no es un polo \( \Rightarrow \) es esencial.

Veamos para \( z_0 = \frac{\pi}{2} + \pi k \) con \( k \in \mathbb{Z} \).
\[
\lim_{z \to \frac{\pi}{2} + \pi k} \frac{\mathrm{sen} z}{\cos z \cdot z^2} = \lim_{z \to \frac{\pi}{2} + \pi k} \frac{\mathrm{sen} z / z^2}{\cos z}
\]
como \( \mathrm{sen} z / z^2 \) está acotado en un entorno de \( \frac{\pi}{2} + \pi k \) y \( \cos z \to 0 \) cuando \( z \to \frac{\pi}{2} + \pi k \)
\[
\Rightarrow \lim_{z \to \frac{\pi}{2} + \pi k} f(z) = \infty
\]
\\
b)
Sin especificar tanto y reutilizando los procedimientos anteriores:

\[
f(z) = \mathrm{sen} \frac{1}{z} \in H(\mathbb{C} \setminus \{0\}) \quad \text{entonces sea } \delta > 0 \Rightarrow f(z) \in H(D'(0, \delta)).
\]

Vemos que \( \lim_{z \to z_0} \mathrm{sen} \frac{1}{z} \)

Sea \( (a_n)_{n=1}^\infty \subset \mathbb{C} \) con \( a_n = \frac{1}{i(n + k)} \) \( \forall n \in \mathbb{N} \cup \{0\} \) con \( k \) suficientemente grande para que \( a_n \subset D'(0, \delta) \). Vemos que \( a_n \xrightarrow{n \to \infty} 0 \)
y
\[
\lim_{n \to \infty} \mathrm{sen} \frac{1}{a_n} = \lim_{n \to \infty} \mathrm{sen}(n+k)i = +\infty
\]

Sin embargo sea \( (b_n)_{n=1}^\infty \subset \mathbb{R} \subset \mathbb{C} \) t.q \( b_n = \frac{1}{2\pi(n+k)} \) con \( k \dots \)
Entonces \( b_n \xrightarrow{n \to \infty} 0 \) pero
\[
\lim_{n \to \infty} \mathrm{sen} \frac{1}{b_n} = \lim_{n \to \infty} \mathrm{sen}(2\pi(n+k)) = 0
\]

\[
\Rightarrow \nexists \lim_{z \to 0} \mathrm{sen} \frac{1}{z}. \quad \text{Es una singularidad \underline{esencial}.}
\]
%ejercicio6
}
\ejercicio{Encuentra el desarrollo de Laurent de $\displaystyle f(z) = \frac{1}{z(1-z)(z-3)}$ en los anillos:
    \begin{enumerate}
        \item $0 < |z| < 1$
        \item $1 < |z| < 2$
        \item $|z| > 2$
    \end{enumerate}}{
a)



aislamos el término de la singularidad en el origen y descomponemos el resto de la función en fracciones parciales.


Escribimos la función como:
\[ f(z) = \frac{1}{z} \cdot \left[ \frac{1}{(1-z)(2-z)} \right] \]
Aplicamos fracciones parciales a la expresión dentro de los corchetes:
\[ \frac{1}{(1-z)(2-z)} = \frac{A}{1-z} + \frac{B}{2-z} \]
Resolviendo para las constantes:
\begin{itemize}
    \item Para $A$: $1 = A(2-z) \implies (\text{evaluando en } z=1) \implies A = 1$.
    \item Para $B$: $1 = B(1-z) \implies (\text{evaluando en } z=2) \implies B = -1$.
\end{itemize}
Por lo tanto:
\[ f(z) = \frac{1}{z} \left( \frac{1}{1-z} - \frac{1}{2-z} \right) \]



En la región indicada, buscamos expansiones que converjan para $|z| < 1$:

\begin{enumerate}
    \item \textbf{Primer término:} Es una serie geométrica estándar:
    \[ \frac{1}{1-z} = \sum_{n=0}^{\infty} z^n \quad \text{válida para } |z| < 1. \]
    
    \item \textbf{Segundo término:} Manipulamos para obtener una forma geométrica:
    \[ \frac{1}{2-z} = \frac{1}{2} \cdot \frac{1}{1 - z/2} = \frac{1}{2} \sum_{n=0}^{\infty} \left( \frac{z}{2} \right)^n = \sum_{n=0}^{\infty} \frac{z^n}{2^{n+1}} \]
    Esta serie converge para $|z/2| < 1$, es decir, $|z| < 2$, lo cual incluye nuestro anillo.
\end{enumerate}



Sustituimos las series en la expresión original:
\[ f(z) = \frac{1}{z} \left[ \sum_{n=0}^{\infty} z^n - \sum_{n=0}^{\infty} \frac{z^n}{2^{n+1}} \right] \]
Agrupamos bajo un mismo sumatorio:
\[ f(z) = \frac{1}{z} \sum_{n=0}^{\infty} \left( 1 - \frac{1}{2^{n+1}} \right) z^n \]
Distribuyendo el factor $1/z$:
\[ f(z) = \sum_{n=0}^{\infty} \left( 1 - \frac{1}{2^{n+1}} \right) z^{n-1} \]


El desarrollo de Laurent para el anillo $0 < |z| < 1$ es:
\[ f(z) = \frac{1}{2z} + \sum_{n=1}^{\infty} \left( 1 - \frac{1}{2^{n+1}} \right) z^{n-1} \]
O reindexando la suma para que comience desde $k=0$:
\[ f(z) = \frac{1}{2z} + \sum_{k=0}^{\infty} \left( 1 - \frac{1}{2^{k+2}} \right) z^k \]
b) Tenemos el anillo: $1 < |z| < 2$, entonces es el anillo $A(0, 1, 2)$.

\noindent Por lo que:
\[
f(z) = \frac{1}{z(1-z)(2-z)} = \frac{1}{z} \left( \frac{1}{1-z} - \frac{1}{2-z} \right)
\]

\noindent Tenemos que como $1 < |z| < 2 \implies 1 > \left| \frac{1}{z} \right| > \frac{1}{2}$. Por lo que:

\[
\frac{1}{1-z} = -\frac{1}{z} \left( \frac{1}{1-\frac{1}{z}} \right) = -\frac{1}{z} \sum_{n=0}^{\infty} \left( \frac{1}{z} \right)^n \quad \text{si } \left| \frac{1}{z} \right| < 1, \text{ es decir, } |z| > 1
\]

\[
\frac{1}{2-z} = \frac{1}{2} \cdot \frac{1}{1-\frac{z}{2}} = \frac{1}{2} \sum_{n=0}^{\infty} \left( \frac{z}{2} \right)^n \quad \text{si } |z| < 2
\]

\noindent Como $z \in A(0, 1, 2)$, entonces:

\[
f(z) = \frac{1}{z} \left( -\frac{1}{z} \sum_{n=0}^{\infty} \left( \frac{1}{z} \right)^n - \frac{1}{2} \sum_{n=0}^{\infty} \left( \frac{z}{2} \right)^n \right) \quad \forall z \in A(0, 1, 2)
\]

\[
f(z) = \frac{1}{z} \left( -\sum_{n=0}^{\infty} \frac{1}{z^{n+1}} - \sum_{n=0}^{\infty} \frac{1}{2^{n+1}} z^n \right)
\]

\[
\iff f(z) = -\sum_{n=0}^{\infty} \frac{1}{z^{n+2}} - \sum_{n=0}^{\infty} \frac{1}{2^{n+1}} z^{n-1}
\]
c)
Partiendo de $f(z) = \frac{1}{z} \left( \frac{1}{1-z} + \frac{1}{z-2} \right)$:
    
    \[ \frac{1}{1-z} = \frac{-1}{z(1 - 1/z)} = -\frac{1}{z} \sum_{n=0}^{\infty} \left( \frac{1}{z} \right)^n = -\sum_{n=0}^{\infty} \frac{1}{z^{n+1}} \]
    
 
    \[ \frac{1}{z-2} = \frac{1}{z(1 - 2/z)} = \frac{1}{z} \sum_{n=0}^{\infty} \left( \frac{2}{z} \right)^n = \sum_{n=0}^{\infty} \frac{2^n}{z^{n+1}} \]



Combinando ambos resultados dentro del factor exterior $1/z$:
\[ f(z) = \frac{1}{z} \left[ \sum_{n=0}^{\infty} \frac{2^n - 1}{z^{n+1}} \right] = \sum_{n=1}^{\infty} \frac{2^n - 1}{z^{n+2}} \]
El desarrollo de Laurent para $|z| > 2$ es:
\[ f(z) = \frac{1}{z^3} + \frac{3}{z^4} + \frac{7}{z^5} + \dots = \sum_{n=1}^{\infty} \frac{2^n - 1}{z^{n+2}} \]

}
\ejercicio{Encuentra el desarrollo de Laurent de $\displaystyle f(z) = \frac{1}{z} + \frac{1}{1-z} + \frac{1}{2-z}$ en los anillos:
    \begin{enumerate}
        \item $0 < |z| < 1$
        \item $0 < |z-1| < 1$
        \item $0 < |z-2| < 1$
    \end{enumerate}
    }{
        Dada $f(z) = \frac{1}{z} + \frac{1}{1-z} + \frac{1}{2-z}$, calculamos el desarrollo de Laurent en tres anillos centrados en sus singularidades.

a) Anillo $0 < |z| < 1$
Desarrollamos en potencias de $z$:
\[ f(z) = \frac{1}{z} + \sum_{n=0}^{\infty} z^n + \sum_{n=0}^{\infty} \frac{z^n}{2^{n+1}} = \frac{1}{z} + \sum_{n=0}^{\infty} \left( 1 + \frac{1}{2^{n+1}} \right) z^n \]

b) Anillo $0 < |z-1| < 1$
Sea $u = z-1$. Expresamos la función en términos de $u$:
\[ f(z) = \frac{1}{1+u} - \frac{1}{u} + \frac{1}{1-u} \]
Desarrollando las series de Taylor para $|u| < 1$:
\[ f(z) = -\frac{1}{u} + \sum_{n=0}^{\infty} (-1)^n u^n + \sum_{n=0}^{\infty} u^n = -\frac{1}{z-1} + \sum_{n=0}^{\infty} ((-1)^n + 1)(z-1)^n \]
Simplificando (solo quedan potencias pares):
\[ f(z) = -\frac{1}{z-1} + \sum_{k=0}^{\infty} 2(z-1)^{2k} \]

c) Anillo $0 < |z-2| < 1$
Sea $v = z-2$. La función queda:
\[ f(z) = \frac{1}{2+v} - \frac{1}{1+v} - \frac{1}{v} \]
Desarrollando para $|v| < 1$:
\[ f(z) = -\frac{1}{v} + \sum_{n=0}^{\infty} \frac{(-1)^n}{2^{n+1}} v^n - \sum_{n=0}^{\infty} (-1)^n v^n \]
Agrupando términos:
\[ f(z) = -\frac{1}{z-2} + \sum_{n=0}^{\infty} (-1)^n \left( \frac{1}{2^{n+1}} - 1 \right) (z-2)^n \]
    }
\ejercicio{Encuentra el desarrollo de Laurent de las funciones siguientes en los anillos que se indican:
    \begin{enumerate}
        \item $\displaystyle \frac{1}{(1-z)(z+2)}$ en $|z|<1$ y en $1<|z|<2$.
        \item $\displaystyle z^3 \cos\left(\frac{1}{z}\right)$ en $|z|>0$.
        \item $\displaystyle \sin\left(\frac{z}{1-z}\right)$ en $|z-1|>0$.
        \item $\displaystyle \frac{1}{z^2+(i-1)z-i}$ en $|z|>1$.
        \item $\displaystyle z e^{-1/z}$ en $|z|>0$.
        \item $\displaystyle \text{Log}\left(1-\frac{2}{z^2}\right)$ en $|z|>\sqrt{2}$.
    \end{enumerate}}{
b)
\noindent Sea $f(z) = z^3 \cos\left(\frac{1}{z}\right)$ en $|z| > 0$, es decir, en $A(0, 0, \infty)$.

\noindent Sabiendo que:
\[
\cos z = \sum_{n=0}^{\infty} (-1)^n \frac{z^{2n}}{(2n)!} \quad \forall z \in \mathbb{C}
\]

\noindent Entonces, como $|z| > 0 \iff z \neq 0 \iff \frac{1}{z} \in \mathbb{C}$, tenemos que:
\[
\cos\left(\frac{1}{z}\right) = \sum_{n=0}^{\infty} \frac{(-1)^n}{(2n)!} \frac{1}{z^{2n}}
\]

\noindent Por lo que:
\[
f(z) = z^3 \cos\left(\frac{1}{z}\right) = z^3 \cdot \sum_{n=0}^{\infty} \frac{(-1)^n}{(2n)!} \frac{1}{z^{2n}} = \sum_{n=0}^{\infty} \frac{(-1)^n}{(2n)!} \frac{1}{z^{2n-3}}
\]

c)
\noindent $f(z) = \text{sen}\left(\frac{z}{1-z}\right)$ en $|z-1| > 0$.

\noindent Tenemos que:
\[
\text{sen}(t) = \sum_{n=0}^{\infty} (-1)^n \frac{t^{2n+1}}{(2n+1)!}
\]

\noindent Además:
\[
\frac{z}{1-z} = -1 + \frac{1}{1-z} \implies \text{sen}\left(\frac{z}{1-z}\right) = \text{sen}\left(-1 + \frac{1}{1-z}\right)
\]

\noindent Entonces, por las propiedades del seno en $\mathbb{C}$ ($\text{sen}(A+B) = \text{sen}A\cos B + \cos A \text{sen}B$):
\[
\text{sen}\left(-1 + \frac{1}{1-z}\right) = \text{sen}(-1)\cos\left(\frac{1}{1-z}\right) + \cos(-1)\text{sen}\left(\frac{1}{1-z}\right)
\]

\noindent Por lo tanto:
\[
f(z) = \sum_{n=0}^{\infty} \text{sen}(-1) \frac{(-1)^n}{(2n)!} (1-z)^{-2n} + \sum_{n=0}^{\infty} \cos(-1) \frac{(-1)^n}{(2n+1)!} (1-z)^{-2n-1}
\]
e)

Partimos del desarrollo de Taylor de la función exponencial $g(w) = e^w$, centrado en $w=0$, el cual tiene un radio de convergencia infinito:
\begin{equation}
e^w = \sum_{n=0}^{\infty} \frac{w^n}{n!} = 1 + w + \frac{w^2}{2!} + \frac{w^3}{3!} + \dots
\end{equation}


Para obtener el término exponencial de nuestra función, realizamos la sustitución $w = -1/z$. Esta sustitución es válida para todo $z$ tal que $z \neq 0$, lo que cumple con la condición de la región $|z| > 0$.
\[ e^{-1/z} = \sum_{n=0}^{\infty} \frac{(-1/z)^n}{n!} = \sum_{n=0}^{\infty} \frac{(-1)^n}{n! z^n} \]


Multiplicamos el desarrollo anterior por el factor $z$ para completar la función original:
\[ f(z) = z \cdot \left( \sum_{n=0}^{\infty} \frac{(-1)^n}{n! z^n} \right) = \sum_{n=0}^{\infty} \frac{(-1)^n}{n! z^{n-1}} \]

Desglosando los términos de la serie de Laurent resultante:
\[ f(z) = z - 1 + \frac{1}{2z} - \frac{1}{6z^2} + \frac{1}{24z^3} - \dots \]

A partir de este desarrollo, observamos que la parte principal de la serie (potencias negativas) tiene infinitos términos. Según la teoría de clasificación de singularidades aisladas, esto demuestra que el punto $z=0$ es una \textbf{singularidad esencial} de la función.
\\
f)
Consideramos la función holomorfa $g(w) = \log(1-w)$ en el disco unidad $|w| < 1$. Es holomorfa puesto que el Logaritmo principal es continuo y holomorfo en $\mathbb{C}$ menos en la semirecta real negativa. Entonces para valores de $w$ en $D(0,1)$ entonces $1-w \in D(1,1)$, haciendo que la composicion de ambas sea holomorfa. 
\begin{equation}
\log(1-w) = -\sum_{n=1}^{\infty} \frac{w^n}{n} = -w - \frac{w^2}{2} - \frac{w^3}{3} - \dots
\end{equation}

Dada la región del ejercicio $|z| > \sqrt{2}$, analizamos el comportamiento del término $w = \frac{2}{z^2}$:
\[ |z| > \sqrt{2} \implies |z|^2 > 2 \implies \left| \frac{2}{z^2} \right| < 1 \]
Puesto que la magnitud de $2/z^2$ es estrictamente menor que 1, estamos dentro del radio de convergencia del desarrollo de Taylor de $g(w)$.


Sustituyendo $w$ por $2/z^2$ en la serie de Taylor, obtenemos la serie de Laurent para $f(z)$:
\[ f(z) = g\left(\frac{2}{z^2}\right) = -\sum_{n=1}^{\infty} \frac{1}{n} \left( \frac{2}{z^2} \right)^n \]
Simplificando la expresión:
\[ f(z) = \sum_{n=1}^{\infty} -\frac{2^n}{n \cdot z^{2n}} \]
El desarrollo explícito comienza con:
\[ f(z) = -\frac{2}{z^2} - \frac{2}{z^4} - \frac{8}{3z^6} - \frac{4}{z^8} - \dots \]
}